\section{总结与展望}

本文围绕“面向简单机械结构建模的自然语言驱动参数化 CAD 生成系统”这一主题，设计并实现了 \texttt{fs-builder} 项目。系统采用大语言模型负责需求理解、强类型 Plan 负责中间约束、确定性模板负责最终脚本生成的架构路线，在自然语言灵活性与 CAD 脚本稳定性之间取得了较好的工程平衡。项目自动化测试全部通过，典型冷拉延模具样例实现了 5 个零件的完整脚本生成，说明系统在限定任务域内具有可行性和实用价值。

本课题的核心价值并不在于实现一个开放域 CAD 智能体，而在于提出并验证了一种适合本科毕业设计场景的工程方法：让大模型只做其擅长的语义理解，让程序化中间层承担约束收缩，让确定性模板负责最终产物落地。

后续工作可从以下几个方向展开：
\begin{enumerate}[leftmargin=2em]
  \item 扩展几何能力，增加孔、沉头孔、倒角、圆角、阵列和螺纹孔等常见特征；
  \item 引入装配约束求解，使系统不仅能表达零件叠放关系，还能进一步处理同轴、平行和间隙等约束；
  \item 接入 Onshape 编译与回读闭环，自动检测 FeatureScript 是否成功编译，并将编译反馈用于纠错；
  \item 增强交互式修订能力，在 Web UI 中加入参数修改、差异对比和增量重生成功能；
  \item 建立评测集，从成功率、字段完整率和脚本有效率等指标系统评估工具质量。
\end{enumerate}

综上，本文完成了一个具备完整闭环、结构清晰、可验证、可展示的毕业设计系统。该系统虽然能力范围有限，但在“自然语言输入 + 受控中间层 + 参数化脚本输出”这一方向上提供了一个具有实践价值的实现样例。
